\begin{proof}
  Define
  \[
    \forall n \in [0, \infty)_{\N}, P(n) \iff \forall \set*{X_{k}}_{k = 1}^{n}, \paren*{\forall k \in [1, n]_{\N},
      \abs*{X_{k}} < \aleph_{0}} \implies \abs*{\prod_{k = 1}^{n} X_{k}} < \aleph_{0}
  \]
  Then
  \begin{itemize}
    \item \textbf{Base case:}
          \[
            \begin{array}{c}
              \prod_{k = 1}^{0} X_{k} = \set*{()}         \\
              \Downarrow                                  \\
              \abs*{\prod_{k = 1}^{0} X_{k}} < \aleph_{0} \\
              \Downarrow                                  \\
              P(0)
            \end{array}
          \]
    \item \textbf{Inductive step:} Suppose
          \begin{itemize}
            \item $P(n)$
            \item $\forall k \in [1, n + 1]_{\N}, \abs*{X_{k}} < \aleph_{0}$
          \end{itemize}
          Since $P(n)$,
          \[
            \abs*{\prod_{k = 1}^{n} X_{k}} < \aleph_{0}
          \]
          By \cref{an1:lem:separated-product-is-similar-to-original-product},
          \[
            \begin{array}{c}
              \prod_{k = 1}^{n + 1} X_{k} \sim \paren*{\prod_{k = 1}^{n} X_{k}} \times X_{n + 1} \\
              \by{\Cref{an1:lem:product-of-two-finite-sets-is-finite}} \Downarrow                \\
              \abs*{\prod_{k = 1}^{n + 1} X_{k}} < \aleph_{0}                                    \\
              \Downarrow                                                                         \\
              P(n + 1)
            \end{array}
          \]
  \end{itemize}
  By \cref{an1:thm:mathematical-induction},
  \[
    \forall n \in [0, \infty)_{\N}, P(n)
  \]
\end{proof}

\begin{theorem} \label{an1:thm:cantors-theorem}
  \[
    \forall X, \nexists f: X \twoheadrightarrow \mathcal{P}(X)
  \]
  This theorem is called \textbf{Cantor's theorem}.
\end{theorem}

\begin{proof}
  Assume
  \[
    \exists X, \exists f: X \twoheadrightarrow \mathcal{P}(X)
  \]
  Define
  \[
    S \coloneqq \setbuilder*{x \in X}{x \notin f(x)} \in \mathcal{P}(X)
  \]
  Then
  \[
    \exists x \in X, f(x) = S
  \]
  This implies
  \begin{itemize}
    \item If $x \in S$, then
          \[
            \begin{array}{c}
              x \notin f(x) = S \\
              \Downarrow        \\
              \bot
            \end{array}
          \]
    \item Otherwise,
          \[
            \begin{array}{c}
              x \in f(x) = S \\
              \Downarrow     \\
              \bot
            \end{array}
          \]
  \end{itemize}
  Therefore,
  \[
    \forall X, \nexists f: X \twoheadrightarrow \mathcal{P}(X)
  \]
\end{proof}
